\chapter{Álgebras de Boole Finitas}

En el capítulo anterior descubrimos que la estructura algebraica pura de Boole encierra en su interior una rigurosa topología de orden. Comprobamos cómo todo par de elementos está sometido a las leyes del supremo ($\vee$) y el ínfimo ($\wedge$) dentro de un retículo acotado, distributivo y complementado.

Esta dualidad álgebra/topología es la clave matemática que nos va a revelar la forma exacta de todas las álgebras de Boole que tienen un número finito de elementos.

\section{Estructura Atómica}

Para comprender la anatomía de un álgebra finita, debemos identificar sus "ladrillos fundamentales", es decir, los elementos indivisibles a partir de los cuales se puede construir todo el conjunto operando con supremos (sumas).

\begin{definicion}{Átomos}{atomos}
Un elemento $x \in \mathbb{B}$ se denomina \textbf{átomo} si es un elemento estrictamente positivo, $x \ne \bot$, y no existe ningún elemento intermedio entre él y el $\bot$. Formalmente:
\[ \mathit{atom}(x) \iff (x > \bot) \wedge \left( \forall y \in \mathbb{B}, \bot < y \le x \implies y = x \right) \]
\end{definicion}

En términos puramente algebraicos, un átomo $x$ es aquel cuyo producto (ínfimo) con cualquier otro elemento $y \in \mathbb{B}$ es el aniquilador total o bien él mismo (absorbente total):
\[ \mathit{atom}(x) \iff (x \ne \bot) \wedge \left( \forall y \in \mathbb{B}, (x \wedge y = \bot) \vee (x \wedge y = x) \right) \]
El conjunto de todos los átomos de un álgebra de Boole se denota como $\mathit{Atom}(\mathbb{B})$.
Una propiedad inmediata es que el ínfimo de dos átomos distintos es siempre nulo:
\[ \forall a, b \in \mathit{Atom}(\mathbb{B}), a \ne b \implies a \wedge b = \bot \]

\begin{definicion}{Hiperátomos (Co-átomos)}{hiperatomos}
De forma dual, un \textbf{hiperátomo} o \textbf{co-átomo} es un elemento estrictamente inferior a $\top$ tal que no existe ningún elemento intermedio entre él y el máximo.
\[ \mathit{hatom}(x) \iff (x < \top) \wedge \left( \forall y \in \mathbb{B}, x \le y < \top \implies y = x \right) \]
\end{definicion}

En cualquier álgebra finita no trivial (donde $\top \ne \bot$), los conjuntos $\mathit{Atom}(\mathbb{B})$ e $\mathit{Hatom}(\mathbb{B})$ nunca están vacíos.

\section{El Teorema de Representación de Stone (Caso Finito)}

Si tomamos el conjunto de todos los átomos $\mathit{Atom}(\mathbb{B})$, podemos generar el conjunto potencia $\wp(\mathit{Atom}(\mathbb{B}))$, es decir, el conjunto de todos sus posibles subconjuntos. Vamos a construir una función $\varphi$ que conecte este álgebra de subconjuntos con el álgebra de Boole original.

\subsection{El Funtor $\varphi$}
Definimos la función $\varphi : \wp(\mathit{Atom}(\mathbb{B})) \to \mathbb{B}$ que toma un subconjunto de átomos $S$ y devuelve el supremo topológico (la suma booleana) de todos ellos:
\[
\varphi(S) = \begin{cases}
\bot & \text{si } S = \emptyset \\
\bigvee_{a \in S} a & \text{si } S \ne \emptyset
\end{cases}
\]

Esta función preserva de forma natural las operaciones del álgebra de subconjuntos hacia las del álgebra de Boole original:
\begin{align*}
\varphi(S_1 \cup S_2) &= \varphi(S_1) \vee \varphi(S_2) \\
\varphi(S_1 \cap S_2) &= \varphi(S_1) \wedge \varphi(S_2) \\
\varphi(\mathit{Atom}(\mathbb{B}) \setminus S) &= \neg \varphi(S)
\end{align*}

\subsection{Isomorfismo y Cardinalidad}
El gran triunfo matemático para las álgebras finitas se resume en el siguiente teorema.

\begin{teorema}{Representación de Álgebras de Boole Finitas}{rep_stone}
Toda álgebra de Boole finita $\mathbb{B}$ es algebraicamente isomorfa al álgebra del conjunto potencia de sus átomos. Es decir, la función $\varphi$ es una biyección perfecta:
\[ \mathbb{B} \simeq \wp(\mathit{Atom}(\mathbb{B})) \]
\end{teorema}
\begin{proof}
Dado que el álgebra es finita, no existen cadenas infinitas descendentes. Todo elemento $x \in \mathbb{B}$ (salvo el $\bot$) acota por arriba a al menos un átomo. Por distributividad y ortogonalidad de los átomos ($a \wedge b = \bot$), cualquier elemento $x$ puede expresarse de forma única como el supremo de los átomos que lo preceden: 
\[ x = \bigvee \{ a \in \mathit{Atom}(\mathbb{B}) \mid a \le x \} \]
Por lo tanto, la función inversa de $\varphi$ está perfectamente definida, demostrando que $\varphi$ es sobreyectiva e inyectiva (biyectiva).
\end{proof}

De este isomorfismo se desprende una consecuencia colosal, formulada previamente de soslayo en nuestro análisis de los conjuntos topológicos.

\begin{teorema}{Cardinalidad de un Álgebra Finita}{cardinalidad}
Si un álgebra de Boole es finita, su número total de elementos debe ser, forzosamente, una potencia de 2.
\[ |\mathbb{B}| = 2^n \quad \text{donde} \quad n = |\mathit{Atom}(\mathbb{B})| \]
\end{teorema}

Cualquier conjunto que no tenga exactamente $2, 4, 8, 16, \ldots$ elementos \textbf{jamás podrá} constituir un álgebra de Boole, sin importar qué operaciones intentemos definir sobre él. Y aún más importante: todas las álgebras de Boole que tengan el mismo número de elementos son exactamente la misma álgebra (son algebraicamente isomorfas). Solo hay \textit{un} álgebra de Boole de 2 elementos, \textit{una} de 4 elementos, \textit{una} de 8 elementos, etc.

Esta uniformidad matemática nos da luz verde para enfocar nuestros esfuerzos en el álgebra bivaluada $\mathbb{B}_2 = \{0, 1\}$ (es decir, la Lógica Proposicional binaria). Sabemos de antemano que cualquier álgebra de cardinal superior no es más que el espacio vectorial n-dimensional $\mathbb{B}_2^n$, un tema que exploraremos a fondo en el siguiente capítulo.
